\section{A localized mixed-trace inverse without pressure coercivity}
\label{vf:section}

\paragraph{Status and consumer.}
This is an author proof, with independent mathematical audit and priority
undetermined. Its input revision is
\texttt{75ca907cad70b7ced4469d8492e7e5dc27cc7444}.
It preserves the displacement-form obstruction of
Section~\ref{va:section}; it does not try to restore that form's positivity.
Instead it solves a specified mixed-trace problem for the \emph{full}
linearized viscous equation on each fixed smooth time interval. The
initial velocity correction is compactly supported. Neither an upper
pressure-Hessian bound nor a spatial Galerkin cutoff is assumed.
This removes finite-interval solvability of this mean-type problem as an
obstruction. It does not supply a singular packet, a uniform inverse at a
singular endpoint, or one common trace for a sequence of different horizons.

Fix $\nu,T>0$. Let $U$ be a real solenoidal field in
$C^\infty([0,T];H^m(\R^3))$ for every integer $m\ge0$.
It can be a known smooth solution or a smooth trial field.
All constants in this section may depend on this \emph{specified} $U$,
$\nu,T$, and the fixed cutoffs. No norm of the as-yet unknown corrected
solution is used to construct the linear inverse.
Write $D_U=\partial_t+U\cdot\nabla$, $M=\nabla U$, and
$\mathcal H=L^2_\sigma(\R^3;\R^3)$. Let
\[
 \mathcal L_U v=\partial_tv-\nu\Delta v+
       \mathbb P\operatorname{div}(U\otimes v+v\otimes U).
\]
Let $E(t,s)$ be its homogeneous forward evolution on $\mathcal H$.
Let $\mathcal T(t,s)$ denote the divergence-preserving transport evolution
$D_U b=M b$. It is the pushforward by the volume-preserving flow of $U$.
Thus $\mathcal T(0,s)=\mathcal T(s,0)^{-1}$ and
\[
 \|\mathcal T(0,s)\|_{2\to2}
 \le \exp\left(\int_0^s\|\nabla U(r)\|_\infty\,dr\right).
\]
This bound concerns the prescribed smooth background on a fixed interval;
it is not offered as a critical continuation estimate.

\begin{lemma}[A positive localized trace operator]\label{vf:localizer}
For real $\chi\in C_c^\infty(\R^3)$ the operator
\[
 A=\operatorname{curl}\,\chi(-\Delta)^{-1}\chi\operatorname{curl}
\]
extends to a bounded nonnegative self-adjoint operator on $\mathcal H$.
It maps every $H^m_\sigma$ into itself and its output is supported in
$\operatorname{supp}\chi$. It is not asserted to be a projection or to
commute with heat evolution.
\end{lemma}
\begin{proof}
Put $C=(-\Delta)^{-1/2}\chi\operatorname{curl}$, so $A=C^*C$ on smooth
tests. The identity
$\chi\operatorname{curl}f=\operatorname{curl}(\chi f)-\nabla\chi\times f$
proves boundedness of $C$ on $L^2$. For the second term, if $g$ has support
in a fixed bounded set, splitting its Fourier transform at radius one gives
\[
 \|(-\Delta)^{-1/2}g\|_2
 \le C(\|g\|_1+\|g\|_2)\le C_\chi\|f\|_2,
 \qquad g=\nabla\chi\times f.
\]
Here $\int_{|\xi|<1}|\xi|^{-2}\,d\xi<\infty$ in three dimensions.
For $C^*g=\operatorname{curl}(\chi(-\Delta)^{-1/2}g)$, the derivative
on the potential is an order-zero Fourier multiplier. The remaining
localized potential is bounded in $L^2$ by the homogeneous Sobolev
inequality $\|(-\Delta)^{-1/2}g\|_6\le C\|g\|_2$ and finite-volume
H\"older. Commuting integer derivatives gives the same bounds with
derivatives of $\chi$, hence $C,C^*$ and $A$ are bounded on $H^m$.
The outer curl proves both the stated support and solenoidality.
Density proves self-adjointness and positivity of the bounded extension.
\end{proof}

\begin{lemma}[Forward evolution and integrated smoothing]\label{vf:smoothing}
The evolutions $E$ and $\mathcal T$ exist on the whole prescribed interval.
For every real $r\ge0$ and $0\le\sigma\le2$,
\[
 \|E(t,s)\|_{H^r\to H^{r+\sigma}}
 \le C_{r,\sigma}(t-s)^{-\sigma/2},\qquad 0\le s<t\le T.
\]
Transport and its inverse are bounded on each $H^r$.
Consequently the operators
\[
 K=\int_0^T\mathcal T(0,s)E(s,0)A\,ds,\qquad
 b_f=\int_0^T\mathcal T(0,s)\int_0^sE(s,r)f(r)\,dr\,ds
\]
satisfy $K:H^r\to H^{r+\sigma}$ for $0\le\sigma<2$ and
$b_f\in H^{r+\sigma}$ whenever $f\in C([0,T];H^r_\sigma)$.
The operator $K:\mathcal H\to\mathcal H$ is compact.
\end{lemma}
\begin{proof}
The linear equation is the Volterra equation with heat evolution and the
actual operator $-\mathbb P\operatorname{div}(U\otimes\cdot+
\cdot\otimes U)$. This operator maps $H^r$ to $H^{r-1}$, and the heat
gain has integrable kernel $C(t-s)^{-1/2}$. Picard iteration on short
intervals and concatenation construct $E$. Differentiated energy estimates,
using bounded spatial derivatives of $U$, give
\[
 \sup_{s\le r\le t}\|v(r)\|_{H^m}^2+
 \nu\int_s^t\|v(r)\|_{H^{m+1}}^2\,dr
 \le C_m\|v(s)\|_{H^m}^2
\]
with the harmless inhomogeneous low-order term incorporated into $C_m$.
On the first half of an interval, this supplies a time with an
$H^{m+1}$ bound $C(t-s)^{-1/2}\|v(s)\|_{H^m}$.
Propagating from that time in $H^{m+1}$ proves one derivative of smoothing.
Splitting the interval into two parts gives two derivatives; interpolation
of Sobolev spaces gives the displayed real-order bounds. Smooth
approximation justifies the argument for $L^2$ inputs.
The flow formula and the chain rule give the transport bounds.
The two integrated smoothing assertions now follow from
$\int_0^T s^{-\sigma/2}\,ds<\infty$ for $\sigma<2$.

We verify compactness without asserting a compact Sobolev embedding on
$\R^3$. Choose $\zeta\in C_c^\infty$ equal to one on
$\operatorname{supp}\chi$. Then $A=\mathbb P\zeta A$.
For each $s>0$, $e^{\nu s\Delta}\mathbb P\zeta$ is Hilbert--Schmidt:
its translation-invariant heat--Leray kernel is in $L^2$, and $\zeta\in L^2$.
In the Volterra series for $E(s,0)\mathbb P\zeta$, each iterated term
is compact. Inductively the integrand is a bounded operator composed with
a compact one for $0<r<s$; the weak singularity at $s$ is integrable,
and the bounded operator norm near zero allows that endpoint to be removed.
Norm convergence of the series, first on short intervals and then by
composition, proves compactness of $E(s,0)\mathbb P\zeta$.
Thus $\mathcal T(0,s)E(s,0)A$ is compact for every $s>0$.
On intervals away from zero it is norm continuous: strong continuity of
transport becomes norm continuity after composition with a compact operator.
The omitted interval $[0,\delta]$ has norm $O(\delta)$.
The integral defining $K$ is therefore a norm limit of compact operators.
\end{proof}

\begin{theorem}[Generic full-PDE mixed-trace solvability]\label{vf:inverse}
There is a locally finite exceptional subset $\mathcal Z\subset(0,\infty)$
such that for every $\lambda\in(0,\infty)\setminus\mathcal Z$ and every
smooth solenoidal $f$ with all spatial Sobolev norms finite, the system
\begin{align*}
 \mathcal L_Uv&=f,& D_U\eta-M\eta&=v,\\
 v(0)&=\lambda A\eta(0),&\eta(T)&=0
\end{align*}
has a unique smooth all-Sobolev solution. Its initial velocity $v(0)$ is
in $C_c^\infty$, with support in $\operatorname{supp}\chi$.
Arbitrarily large admissible positive $\lambda$ exist.
No pressure-curvature hypothesis is required.
\end{theorem}
\begin{proof}
Variation of constants gives
\[
 v(s)=E(s,0)\lambda A\eta_0+v_f(s),\qquad
 v_f(s)=\int_0^sE(s,r)f(r)\,dr.
\]
Solving the displacement equation and pulling its terminal value back to
time zero shows that its sole compatibility equation is
\begin{equation}\label{vf:schur}
 (I+\lambda K)\eta_0=-b_f.
\end{equation}
All propagators act on the original solenoidal whole-space Hilbert space.
In particular this formula retains the mean pressure and all exterior modes.

Here is a finite-dimensional proof of the required compact-operator
alternative. Complexify $\mathcal H$. On any disk $|z|\le R$, approximate
$K$ in operator norm by a finite-rank $K_N$ with
$\|K-K_N\|<1/(4R)$. Put $B_z=(I+z(K-K_N))^{-1}$, defined by its
Neumann series on the larger disk $|z|<4R$.
Let $W=\operatorname{ran}K_N$. The factorization
\[
 I+zK=(I+z(K-K_N))(I+zB_zK_N)
\]
and finite-rank linear algebra show that invertibility is equivalent to
nonvanishing of
\[
 D(z)=\det\bigl(I_W+zK_NB_z|_W\bigr).
\]
This analytic function has $D(0)=1$, so has only finitely many zeros on
$[0,R]$. The same argument on each larger disk proves that the positive
exceptional set is locally finite and proves both existence and uniqueness
off that set. Since all operators are real, uniqueness makes the inverse
real on real inputs.

Equation~\eqref{vf:schur} initially gives $\eta_0\in L^2$.
The relation $\eta_0=-b_f-\lambda K\eta_0$ and the gain of one
spatial derivative from Lemma~\ref{vf:smoothing} inductively give every
$H^m$ norm. The forward and transport equations give the claimed solution
regularity. Lemma~\ref{vf:localizer} then gives the compact smooth initial
velocity. Conversely every solution must satisfy \eqref{vf:schur}, so
there are no additional compatibility conditions or hidden initial traces.
\end{proof}

\begin{proposition}[Heat background: every positive penalty works]\label{vf:heat}
When $U=0$, set $G=\int_0^T e^{\nu s\Delta}\,ds$.
Then $K=GA$, and for every $\lambda\ge0$,
\[
 (I+\lambda GA)^{-1}
 =I-\lambda G A^{1/2}
   (I+\lambda A^{1/2}GA^{1/2})^{-1}A^{1/2},
 \qquad
 \|(I+\lambda GA)^{-1}\|\le1+\lambda T\|A\|.
\]
\end{proposition}
\begin{proof}
Both $G$ and $A$ are bounded nonnegative self-adjoint operators.
Hence $I+\lambda A^{1/2}GA^{1/2}$ has inverse norm at most one.
Multiplication verifies the stated formula; $\|G\|\le T$ proves the bound.
Commutation of $G$ and $A$ is neither assumed nor needed.
\end{proof}

\paragraph{What is and is not repaired.}
The inverse is a non-self-adjoint forward-propagator construction, not a
positive displacement action. It therefore does not contradict
Theorem~\ref{va:classification}, even for the arbitrarily small genuine
NS backgrounds used there. Its trace is a \emph{velocity} condition at
time zero; the displacement is transported using the full background.
No compact support is imposed on positive-time velocity or pressure.
For countably many fixed smooth histories, a common positive $\lambda$
can avoid their countably many exceptional sets. This does \emph{not}
make their initial velocities equal or their inverse norms uniformly bounded.
The first outstanding use in UE1 remains quantitative inversion and
Cauchy-trace coherence on a concentrating sequence. No forward amplification,
blowup-preserving nonlinear correction, or arbitrary-data RF-$q$ bound is
obtained merely from Theorem~\ref{vf:inverse}.
